\documentclass[12pt]{article}
\usepackage[utf8]{inputenc}
\usepackage[T1]{fontenc}
\usepackage[spanish]{babel}
\usepackage{geometry}
\usepackage{amsmath,amssymb}
\usepackage{enumitem}

\geometry{margin=2.5cm}

\begin{document}

\begin{center}
{\Large Transcripción cruda}\\[2mm]
{\large 21May26}
\end{center}

\section*{Fotos 18--20}

\textit{Más ejercicios:} 22, 23, 24, 25, 26, 27, 29.

\textit{Proposición.} Sea $n\in\mathbb N$ fijo. Sean $x,y\in\mathbb R^n$,
\[
x=(x_1,\dots,x_n),\qquad y=(y_1,\dots,y_n).
\]
Entonces
\[
\max_{1\le i\le n}|x_i-y_i|
\le
\left(\sum_{i=1}^n (x_i-y_i)^2\right)^{1/2}
\le
\sqrt n\,\max_{1\le i\le n}|x_i-y_i|.
\]
La primera es la llamada nueva distancia y la segunda es la distancia euclidiana.

\textit{Dem.} Si
\[
m=\max_{1\le i\le n}|x_i-y_i|,
\]
entonces
\[
m^2=\max_{1\le i\le n}|x_i-y_i|^2
\le \sum_{i=1}^n (x_i-y_i)^2
\le n\,m^2.
\]
Extrayendo raíz cuadrada se obtiene la prueba.

\textit{Teorema.} Sea $(X,d_1)$ un espacio métrico. Sea $d_2:X\times X\to[0,\infty)$.
Supongamos que existen constantes $c$ y $C$ positivas, con $c<C$, tales que
\[
c\,d_1(x,y)\le d_2(x,y)\le C\,d_1(x,y),\qquad \forall x,y\in X.
\]
Entonces $d_1$ y $d_2$ son equivalentes.

\textit{Dem.} Requerimos varias etapas.

\textit{Lema.} Si $(X,d)$ es espacio métrico, sea $A\subset X$ un conjunto abierto. Entonces $A$ es la unión de vecindades, es decir, es la unión de conjuntos de la forma $N_r(q)$.

\textit{Dem. del lema.} Sea $q\in A$. Existe $r_q>0$ tal que $N_{r_q}(q)\subset A$.
Luego
\[
A\subset \bigcup_{q\in A} N_{r_q}(q)\subset A,
\]
así que
\[
A=\bigcup_{q\in A} N_{r_q}(q).
\]

\textit{Dem. del teorema.} Sea $A\subset X$ abierto en $(X,d_1)$. Probamos que $A$ es abierto en $(X,d_2)$.
Sea $a\in A$. Entonces existe $r_a>0$ tal que $N_{r_a}^{d_1}(a)\subset A$.
Como
\[
N_{r_a/C}^{d_2}(a)\subset N_{r_a}^{d_1}(a),
\]
se obtiene
\[
N_{r_a/C}^{d_2}(a)\subset A.
\]
Recíprocamente, si $A$ es abierto en la métrica $d_2$, entonces existe $r_a>0$ tal que
\[
N_{r_a}^{d_2}(a)\subset A.
\]
Como
\[
N_{r_a/C}^{d_1}(a)\subset N_{r_a}^{d_2}(a),
\]
se concluye que $A$ es abierto en $(X,d_1)$.

\textit{Corolario.} Si $F\subset X$ es cerrado y $K\subset X$ es compacto, entonces $F\cap K$ es compacto.

\textit{Dem.} $F\cap K$ es cerrado porque es intersección de cerrado con compacto cerrado. Además $F\cap K\subset K$. Por el teorema anterior, $F\cap K$ es compacto.

\textit{Teorema.} Sea $\mathcal K=\{K_\alpha\}_{\alpha\in\Lambda}$ una familia de subconjuntos compactos de $X$.
Supongamos que toda subfamilia finita de $\mathcal K$ tiene intersección no vacía. Entonces
\[
\bigcap_{\alpha\in\Lambda}K_\alpha\neq\varnothing.
\]

\textit{Dem.} Si
\[
\bigcap_{\alpha\in\Lambda}K_\alpha=\varnothing,
\]
entonces, para algún $\alpha_0\in\Lambda$, se tiene
\[
K_{\alpha_0}\cap\bigcap_{\alpha\in\Lambda\setminus\{\alpha_0\}}K_\alpha=\varnothing.
\]
Sea $\alpha_0\in\Lambda$ fijo. Entonces
\[
\varnothing=K_{\alpha_0}\cap\bigcap_{\alpha\in\Lambda\setminus\{\alpha_0\}}K_\alpha.
\]
Luego
\[
K_{\alpha_0}\subset \bigcup_{\alpha\in\Lambda\setminus\{\alpha_0\}}K_\alpha^c.
\]
Como $K_\alpha^c$ es abierto, la familia $\{K_\alpha^c\}_{\alpha\in\Lambda\setminus\{\alpha_0\}}$ es una cubierta abierta de $K_{\alpha_0}$.
Por compacidad de $K_{\alpha_0}$, existen $\alpha_1,\dots,\alpha_m\in\Lambda$ tales que
\[
K_{\alpha_0}\subset K_{\alpha_1}^c\cup\cdots\cup K_{\alpha_m}^c.
\]
Por tanto,
\[
K_{\alpha_0}\cap K_{\alpha_1}\cap\cdots\cap K_{\alpha_m}=\varnothing,
\]
contradicción.

\section*{Fotos 1--4}

\textit{Proposición.} Sea $(X,d_1)$ un espacio métrico. Sea
\[
d_2:X\times X\to[0,\infty),\qquad
d_2(x,y)=\frac{d_1(x,y)}{1+d_1(x,y)}.
\]
Entonces $d_2$ es métrica.

\textit{Dem.} Positividad:
\[
d_2(x,y)\ge 0,\qquad d_2(x,y)=0 \iff d_1(x,y)=0 \iff x=y.
\]
Simetría:
\[
d_2(x,y)=d_2(y,x).
\]
Falta la desigualdad del triángulo:
\[
d_2(x,y)\le d_2(x,z)+d_2(z,y),\qquad \forall x,y,z\in X.
\]
Es decir, hay que probar
\[
\frac{d_1(x,y)}{1+d_1(x,y)}
\le
\frac{d_1(x,z)}{1+d_1(x,z)}
+\frac{d_1(z,y)}{1+d_1(z,y)}.
\]

\textit{Introducimos tres funciones}
\[
f,g,h:\mathbb R_+\to\mathbb R_+,
\]
\[
f(a)=\frac a{1+a}=1-\frac1{1+a}.
\]
Sea $b>0$ fijo.
\[
g(c)=\frac{b+c}{1+b+c},
\qquad
h(c)=\frac b{1+b}+\frac c{1+c}.
\]
Demostraremos que
\[
g(c)\le h(c)\qquad \forall c\ge 0.
\]
Si $0\le a\le b+c$, entonces
\[
\frac a{1+a}=f(a)\le f(b+c)=g(c)\le h(c)=\frac b{1+b}+\frac c{1+c}.
\]

\textit{Demostración.} Consideramos
\[
g(0)=\frac b{1+b},\qquad h(0)=\frac b{1+b}.
\]
Además
\[
g'(c)=\frac1{(1+b+c)^2},\qquad h'(c)=\frac1{(1+c)^2}.
\]
Como $1+b+c>1+c$, se tiene
\[
g'(c)<h'(c).
\]
Por el teorema fundamental del cálculo,
\[
g(c)=g(0)+\int_0^c g'(x)\,dx
\le h(0)+\int_0^c h'(x)\,dx
h(c).
\]
Así se obtiene la desigualdad requerida.

\section*{Fotos 5--8}

\textit{Proposición.} Sea $(X,d_1)$ espacio métrico. Sea
\[
d_2(x,y)=\frac{d_1(x,y)}{1+d_1(x,y)}.
\]
Entonces $d_2$ es métrica.

\textit{Dem.} Ya vimos positividad y simetría. La desigualdad del triángulo se reduce a la desigualdad anterior para las funciones $f,g,h$.

\textit{Teorema.} Sea $(X,d_1)$ espacio métrico y sea $K\subset X$ un conjunto compacto. Entonces $K$ es cerrado.

\textit{Dem.} Probamos que $K^c$ es abierto.
Sea $p\in K^c$.
Para cada $q\in K$, sea
\[
V_q=\frac12 d(p,q)>0.
\]
Afirmación:
\[
N_{V_q}(q)\cap N_{V_q}(p)=\varnothing.
\]
Es claro que
\[
K\subset \bigcup_{q\in K}N_{V_q}(q).
\]
Por compacidad de $K$, existen $q_1,\dots,q_n\in K$ tales que
\[
K\subset N_{V_{q_1}}(q_1)\cup\cdots\cup N_{V_{q_n}}(q_n).
\]
Como
\[
p\in\bigcap_{i=1}^n N_{V_{q_i}}(p),
\]
se demuestra que
\[
\bigcap_{i=1}^n N_{V_{q_i}}(p)\cap K=\varnothing.
\]
Por tanto $K^c$ es abierto.

\section*{Fotos 9--12}

\textit{Teorema.} Sea $(X,d_1)$ y sea $d_2$ otra métrica sobre $X$.
Supongamos que existen constantes $c$ y $C$ positivas tales que
\[
c\,d_1(x,y)\le d_2(x,y)\le C\,d_1(x,y),\qquad \forall x,y\in X.
\]
Entonces $d_1$ y $d_2$ son equivalentes.

\textit{Lema.} Si $(X,d)$ es espacio métrico, sea $A\subset X$ un conjunto abierto. Entonces $A$ es la unión de vecindades; es decir, la unión de conjuntos de la forma $N_r(q)$.

\textit{Dem. del lema.} Si $q\in A$, existe $r_q>0$ tal que $N_{r_q}(q)\subset A$.
Luego
\[
A=\bigcup_{q\in A}N_{r_q}(q).
\]

\textit{Dem. del teorema.} Sea $A\subset X$ abierto en $(X,d_1)$. Probamos que $A$ es abierto en $(X,d_2)$.
Sea $a\in A$ y sea $r_a>0$ tal que
\[
N_{r_a}^{d_1}(a)\subset A.
\]
Como
\[
N_{r_a/C}^{d_2}(a)\subset N_{r_a}^{d_1}(a),
\]
se sigue que $A$ es abierto en $(X,d_2)$.
Recíprocamente, si $A$ es abierto en $(X,d_2)$, entonces existe $r_a>0$ tal que
\[
N_{r_a}^{d_2}(a)\subset A,
\]
y como
\[
N_{r_a/C}^{d_1}(a)\subset N_{r_a}^{d_2}(a),
\]
se obtiene que $A$ es abierto en $(X,d_1)$.

\section*{Fotos 13--17}

\textit{Definición.} Sea $(X,d)$ un espacio métrico. Sea $D\subset X$. Decimos que $D$ es denso si
\[
\overline D=X.
\]
También se escribe
\[
D\cup D'= \overline D = X.
\]
\textit{b)} Un espacio métrico $(X,d)$ es separable si existe $D\subset X$ tal que $D$ es denso y es numerable.

\textit{Ejemplos.}
\[
E_1=\left\{\frac1n:n\in\mathbb N\right\}\cup\{0\},\qquad E_1'=\{0\},
\]
\[
E_2=\left\{2+\frac1n:n\in\mathbb N\right\}\cup\{2\},\qquad E_2'=\{2\},
\]
\[
E_3=\left\{10+\frac1n:n\in\mathbb N\right\}\cup\{10\},\qquad E_3'=\{10\}.
\]
Y también
\[
E_1\cup E_2\cup E_3.
\]

\[
F_j=\left\{\frac1j+\frac1n:n=1,2,\dots\right\}\cup\left\{\frac1j\right\},
\qquad
F_j'=\left\{\frac1j\right\}.
\]
Finalmente,
\[
\bigcup_{j=1}^\infty F_j
\]
da un ejemplo de unión numerable.

\end{document}
